% --- ПРЕАМБУЛА ДЛЯ ПОДДЕРЖКИ РУССКОГО ЯЗЫКА И МАТЕМАТИКИ ---
\usepackage[T2A]{fontenc}
\usepackage[utf8]{inputenc}
\usepackage[russian]{babel}
\usepackage{amsmath}
\usepackage{geometry}
\usepackage{amssymb}
\usepackage{graphicx}
\usepackage{float}
\usepackage{import}
\geometry{a4paper, top=2cm, bottom=2cm, left=2cm, right=2cm}
\newcommand{\sidecomment}[1]{\marginpar{\raggedright\scriptsize\itshape #1}}
\newcommand{\iu}{\mathrm{i}}

\newcommand{\Def}{%
    \text{Def}%
    \hspace{-0.85em}% Сдвиг назад ровно на середину (подберите под шрифт)
    \raisebox{0.1ex}{/}% Черта
    \hspace{1em}% Возврат к концу слова + небольшой отступ
}

\renewcommand{\Re}{\operatorname{Re}}
\renewcommand{\Im}{\operatorname{Im}}

\ifstandalone
    \newcommand{\lecturebreak}{} % В отдельном файле лекции разрыв не нужен
\else
    \newcommand{\lecturebreak}{\clearpage} % В полной книге - начинаем с новой страницы
\fi

\newcounter{lecturenumber} % Создаем счетчик для нумерации лекций
\newcommand{\lecture}[2]{
    \lecturebreak
    \refstepcounter{lecturenumber} % Увеличиваем счетчик на 1
    \par\vspace{2em}\noindent % Отступ сверху и убираем абзацный отступ
    {\Large\bfseries % Делаем заголовок крупным и жирным
    Лекция \thelecturenumber: #1 % Выводим "Лекция N: Тема"
    \hfill % Растягиваем пробел, чтобы дата была справа
    \large\textit{#2}} % Выводим дату курсивом
    \par\vspace{1em} % Небольшой отступ снизу
    \addcontentsline{toc}{section}{Лекция \thelecturenumber: #1} % Добавляем в оглавление
    \par
}

\newcommand{\TwoCols}[2]{%
  \par\noindent % Начать с новой строки без отступа
  \begin{minipage}[t]{0.48\textwidth}
    #1
  \end{minipage}% <-- Важный '%' для удаления пробела
  \hfill % Растягивающийся пробел между колонками
  \begin{minipage}[t]{0.48\textwidth}
    #2
  \end{minipage}% <-- Важный '%' для удаления пробела
  \par%
}

\pagestyle{plain} % Включаем нумерацию страниц\